\documentclass[10pt,a4paper]{amsart}
\usepackage[margin=19mm]{geometry}
\usepackage{amsmath,amssymb,mathtools}
\usepackage[scr=rsfs,cal=euler]{mathalfa}
\usepackage{xfrac}
\usepackage[unicode,hidelinks,bookmarks=false]{hyperref}
\newcommand{\ie}{i.e.\ }
\newcommand{\cf}{cf.\ }
\newcommand{\resp}{respectively\ }
\newcommand{\Subsection}{\subsection}
\providecommand{\nobreakdash}{\nobreak\hbox{-}}
\setlength{\emergencystretch}{2em}
\multlinegap0pt
\usepackage[normalem]{ulem}
\usepackage{amssymb}
\usepackage{enumerate}
\usepackage{natbib}
\usepackage{float}
\usepackage[shortlabels]{enumitem}
\usepackage{xcolor}
\newtheorem{theorem}{Theorem}
\title{Oscillatory Regimes in a Game-Theoretic Model for Mosquito Population Dynamics under Breeding Site Control}
\author{Mohammad Rubayet Rahman, Chanaka Kottegoda,, Lucas M. Stolerman}
\date{Source: arXiv:2601.05222v1 (CC BY 4.0)}
\begin{document}
\maketitle

\noindent\emph{Source and licence.} Oscillatory Regimes in a Game-Theoretic Model for Mosquito Population Dynamics under Breeding Site Control. Version arXiv:2601.05222v1 (CC BY 4.0), distributed by arXiv under the Creative Commons Attribution 4.0 International licence, http://creativecommons.org/licenses/by/4.0/, as recorded in the arXiv licence metadata for this exact version and shown on the abstract page https://arxiv.org/abs/2601.05222v1. Full licence notice: \texttt{licenses/arxiv-2601.05222-CC-BY-4.0.txt}.\par\smallskip

\noindent\textit{Editorial scope.} Counted unit: the article's theorem Model (source0.tex lines 581--594). The article's complete original proof of it is delivered with it (source0.tex lines 987--1267, one \texttt{proof} environment of 281 lines, which the article places away from the statement). No mathematics is written, repaired or re-derived: the statement and the proof are the article's own lines, in the article's own notation, and no retained line is substituted or re-flowed.  Retained uncounted as the source-local context the proof needs: lines 185--193.  Nothing was omitted from any retained span, so the package carries no omission record.  No retained line contains a citation, so the wrapper carries no bibliography and no bibliography entry.  No retained line contains a figure environment, an image-inclusion command or drawing code, so no image is delivered and the package declares no asset dependency.  Editorial formatting only, in addition: an AMS \texttt{amsart} preamble that restates the article's own declarations and macros used by the retained lines, namely the environments \texttt{theorem} on separate counters, together with the standard packages those lines need.  The retained environments print in this document as Theorem 1 in order of appearance, whereas the article prints the same items with the numbering of its own full document.  The complete native source of the article as distributed in the arXiv e-print for this exact version is bundled unchanged as \texttt{sources/192329/source0.tex}.
The model equations are given by the following system: 
\begin{equation}
\begin{aligned}
\frac{dL_v}{dt} &= rbA_v \left(1 - \frac{L_v}{K_v(w)} \right) - (\nu_L + \mu_L)L_v, \\
\frac{dA_v}{dt} &= \nu_L L_v - \mu_A A_v, \\
\frac{dw}{dt}   &= k w(1 - w)\left[-r_c + r_d m A_v\right].
\end{aligned}
\label{eq:behavioral_entomological_system}
\end{equation}

\begin{theorem}\label{thm: Game-theoretic Model}
The local stability of $E_{01}$ -- $E_{04}$ for System~\eqref{eq:behavioral_entomological_system} is characterized as follows:
\begin{enumerate}
 \item \( E_{01} \) is LAS if \( N < 1 \) and unstable if \( N > 1 \).
    \item \( E_{02} \) is unstable.
    
    \item $E_{03}$ is LAS if  $\frac{r_c}{r_d} > \alpha K_{v_{\max}}\left(1 - \frac{1}{N}\right)$ and unstable if $\frac{r_c}{r_d} < \alpha K_{v_{\max}}\left(1 - \frac{1}{N}\right).$

    \item $E_{04}$ is LAS if  $\frac{r_c}{r_d} < \alpha K_{v_{\min}}\left(1 - \frac{1}{N}\right)$ and unstable if $\frac{r_c}{r_d} > \alpha K_{v_{\min}}\left(1 - \frac{1}{N}\right).$

    
    
\end{enumerate}
\end{theorem}

\begin{proof}[\textbf{Proof of Theorem \ref{thm: Game-theoretic Model}}]

 We proceed with the analysis of each steady state separately, proving the claims from items 1--4. 

 
\begin{itemize}


\item \underline{$E_{01}$:} We evaluate the Jacobian matrix at the steady state $E_{01}$. The Jacobian matrix at $E_{01}$ is given by

\[
J = \begin{bmatrix}
- (\nu_L + \mu_L) & r b & 0 \\
\nu_L & -\mu_A & 0 \\
0 & 0 & -k r_c
\end{bmatrix}
\]

and the characteristic equation is given by \(\det(J - \lambda I) = 0\) where 

\[
J - \lambda I = \begin{bmatrix}
- (\nu_L + \mu_L) - \lambda & r b & 0 \\
\nu_L & -\mu_A - \lambda & 0 \\
0 & 0 & -k r_c - \lambda
\end{bmatrix}.
\]

The determinant of this matrix can be calculated by expanding along the third row, i.e.,

$$ 
\begin{aligned}
\operatorname{det}(J-\lambda I) & =(-k r_c-\lambda)  \operatorname{det}\left[\begin{array}{cc}
-\left(\nu_L+\mu_L\right)-\lambda & r b \\
\nu_L & -\mu_A-\lambda
\end{array}\right] \\
& =(-k r_c-\lambda)\left(-\left(\nu_L+\mu_L\right)-\lambda\right)\left(-\mu_A-\lambda\right)-(r b)\left(\nu_L\right) \\
& =(-k r_c-\lambda)\left(\lambda^2+\left(\left(\nu_L+\mu_L\right)+\mu_A\right) \lambda+\left(\nu_L+\mu_L\right) \mu_A-r b \nu_L\right).
\end{aligned}
$$


The eigenvalues are \(\lambda_1 = -k r_c\) and 

$$
\lambda_{2,3} = \frac{-\left(\nu_L + \mu_L + \mu_A\right) \pm \sqrt{\left(\nu_L + \mu_L + \mu_A\right)^2 - 4\left((\nu_L + \mu_L) \mu_A - r b \nu_L\right)}}{2}.
$$
Since $k>0$, the eigenvalue \(\lambda_1\) will be negative if and only if \(-k r_c < 0 \Rightarrow r_c > 0\). Moreover, since  $(\nu_L + \mu_L + \mu_A) > 0 $,  the other two eigenvalues \(\lambda_{2,3}\) will have negative real parts if and only if
$$
\begin{aligned}
& \left(\nu_L+\mu_L\right) \mu_A-r b \nu_L>0 
& \Leftrightarrow\quad {N}<1.
\end{aligned}
$$
Since $r_c > 0$ is also our model assumption, the analysis confirms that $E_{01}$ is LAS, establishing the stability conditions stated in item 2. On the other hand, $E_{01}$ loses stability if $N>1$.


 \item \underline{$E_{02}$:}  We evaluate the Jacobian matrix at $E_{02}$ and analyze the associated eigenvalues. Then we determine the conditions under which the steady state is LAS. The Jacobian matrix at $E_{02}$ is given by


\[
J = \begin{bmatrix}
-(\nu_L + \mu_L) & r b & 0 \\
\nu_L & -\mu_A & 0 \\
0 & 0 & k r_c
\end{bmatrix}.
\]

Now the characteristic equation is given by \(\det(J - \lambda I) = 0\) where 

\[
J - \lambda I = \begin{bmatrix}
-(\nu_L + \mu_L) - \lambda & r b & 0 \\
\nu_L & -\mu_A - \lambda & 0 \\
0 & 0 & k r_c - \lambda
\end{bmatrix}.
\]

The determinant of this matrix can be calculated by expanding along the third row, i.e., 

$$
\begin{aligned}
\operatorname{det}(J-\lambda I) & =\left(k r_c-\lambda\right) \operatorname{det}\left[\begin{array}{cc}
-\left(\nu_L+\mu_L\right)-\lambda & r b \\
\nu_L & -\mu_A-\lambda
\end{array}\right] \\
& =\left(k r_c-\lambda\right)\left(\lambda^2+\left[\left(\nu_L+\mu_L\right)+\mu_A\right] \lambda+\left[\left(\nu_L+\mu_L\right) \mu_A-r b \nu_L\right]\right).
\end{aligned}
$$
The first eigenvalue \(\lambda_1 = k r_c\) and the other two eigenvalues 
$$
\lambda_{2,3} = \frac{-\left((\nu_L + \mu_L) + \mu_A\right) \pm \sqrt{\left((\nu_L + \mu_L) + \mu_A\right)^2 - 4\left((\nu_L + \mu_L) \mu_A - r b \nu_L\right)}}{2}.$$
 
Since $k>0$, so the eigenvalue \(\lambda_1\) will be negative if and only if \(k r_c < 0 \Leftrightarrow  r_c < 0
\) and \(\lambda_{2,3}\) will have negative real parts if and only if $$
\left(\left(\nu_L+\mu_L\right) \mu_A-r b \nu_L\right)>0 
 \Leftrightarrow \quad {N}<1.
$$
Here $r_c < 0$ is not feasible since the perceived risk is assumed to be a positive real number. Hence $E_{02}$ is always unstable.




 \item \underline{$E_{03}$:} Next, the Jacobian matrix at $E_{03}$ is given by

$$
J=\left[\begin{array}{ccc}
-\left(\frac{r b \nu_L}{\mu_A}\left(1-\frac{1}{{N}}\right)+\nu_L+\mu_L\right) & \frac{r b}{{N}} & -r b \frac{\nu_L}{\mu_A}\left(1-\frac{1}{{N}}\right)^2\left(K_{v_{\max }}-K_{v_{\text {min }}}\right) \\
\nu_L & -\mu_A & 0 \\
0 & 0 & k\left(-r_c+r_d m\left(\frac{\nu_L}{\mu_A} K_{v_{\text {max }}}\left(1-\frac{1}{{N}}\right)\right)\right.
\end{array}\right]
$$

and the characteristic equation is given by \(\det(J - \lambda I) = 0\) where

\[
%\small{ 
J-\lambda I = \left[
\begin{array}{ccc}
-\left(\frac{r b \nu_L}{\mu_A}\left(1 - \frac{1}{{N}}\right) + \nu_L + \mu_L\right) - \lambda 
& \frac{r b}{{N}} 
& -r b \frac{\nu_L}{\mu_A}\left(1 - \frac{1}{{N}}\right)^2\left(K_{v_{\max}} - K_{v_{\min}}\right) \\
\nu_L 
& -\mu_A - \lambda 
& 0 \\
0 
& 0 
& k\left(-r_c + r_d  m \left(\frac{\nu_L}{\mu_A} K_{v_{\max}}\left(1 - \frac{1}{{N}}\right)\right)\right) - \lambda
\end{array}
\right].
%}
\]

Now the determinant of this matrix can be calculated by expanding along the third row, i.e.,

$$
\begin{aligned}
& \operatorname{det}(J-\lambda I)=\left[k\left(-r_c+r_d m\left(\frac{\nu_L}{\mu_A} K_{v_{\max }}\left(1-\frac{1}{{N}}\right)\right)\right)-\lambda\right] \\
& \operatorname{det}\left[\begin{array}{cc}
-\left(\frac{r b \nu_L}{\mu_A}\left(1-\frac{1}{{N}}\right)+\nu_L+\mu_L\right)-\lambda & \frac{r b}{{N}} \\
\nu_L & -\mu_A-\lambda
\end{array}\right]
\end{aligned}
$$

$$
\begin{aligned}
&=\left[k\left(-r_c+r_d m\left(\frac{\nu_L}{\mu_A} K_{v_{\max }}\left(1-\frac{1}{{N}}\right)\right)\right)-\lambda\right]\\
& {\left[\left(\frac{r b \nu_L}{\mu_A}\left(1-\frac{1}{{N}}\right)+\nu_L+\mu_L\right)+\lambda\right]\left[\mu_A+\lambda\right]-\left[\frac{r b}{{N}}\right]\left[\nu_L\right] }
\end{aligned}
$$

$$
\begin{aligned}
= & {\left[k\left(-r_c+r_d m\left(\frac{\nu_L}{\mu_A} K_{v_{\max }}\left(1-\frac{1}{{N}}\right)\right)\right)-\lambda\right] } \\
& {\left[\left(\frac{r b \nu_L}{\mu_A}\left(1-\frac{1}{{N}}\right)+\nu_L+\mu_L\right) \mu_A+\left(\frac{r b \nu_L}{\mu_A}\left(1-\frac{1}{{N}}\right)+\nu_L+\mu_L\right) \lambda+\mu_A \lambda+\lambda^2-\frac{r b \nu_L}{{N}}\right] }
\end{aligned}
$$

$$
\begin{aligned} 
& =\left[k\left(-r_c+r_d  m\left(\frac{\nu_L}{\mu_A} K_{v_{\max }}\left(1-\frac{1}{{N}}\right)\right)\right)-\lambda\right]  \\
& {\left[\lambda^2+\left(\left(\frac{r b \nu_L}{\mu_A}\left(1-\frac{1}{{N}}\right)+\nu_L+\mu_L\right)+\mu_A\right) \lambda+\left(\left(\frac{r b \nu_L}{\mu_A}\left(1-\frac{1}{{N}}\right)+\nu_L+\mu_L\right) \mu_A-\frac{r b \nu_L}{{N}}\right)\right]}.
\end{aligned}
$$

Hence, we have eigenvalues $\lambda_1=k\left(-r_c+r_d m\left(\frac{\nu_L}{\mu_A} K_{v_{\max }}\left(1-\frac{1}{{N}}\right)\right)\right)
$ which is true if $N>1$  and

$
\begin{aligned}
& \lambda_{2,3}=\frac{-\left(\left(\frac{r b \nu_L}{\mu_A}\left(1-\frac{1}{{N}}\right)+\nu_L+\mu_L\right)+\mu_A\right)}{2} \pm \\
& \frac{\sqrt{\left(\left(\frac{r b \nu_L}{\mu_A}\left(1-\frac{1}{{N}}\right)+\nu_L+\mu_L\right)+\mu_A\right)^2-4\left(\left(\frac{r b \nu_L}{\mu_A}\left(1-\frac{1}{{N}}\right)+\nu_L+\mu_L\right) \mu_A-\frac{r b \nu_L}{{N}}\right)}}{2}.
\end{aligned}
$

Here, eigenvalue \(\lambda_1\) will be negative if and only if $$
\begin{aligned}
& k\left[-r_c+r_d m\left(\frac{\nu_L}{\mu_A} K_{v_{\max }}\left(1-\frac{1}{{N}}\right)\right)\right]<0 \\
& \Rightarrow \frac{r_c}{r_d}>\frac{m \nu_L}{\mu_A} K_{v_{\max }}\left(1-\frac{1}{{N}}\right) \\
& \Rightarrow \frac{r_c}{r_d}>\alpha K_{v_{\max }}\left(1-\frac{1}{{N}}\right) \text { where } \alpha=\frac{m \nu_L}{\mu_A}
\end{aligned}
$$

and the real parts of the eigenvalues \(\lambda_{2, 3}\) will be negative if and only if
$$
\left(\frac{r b \nu_L}{\mu_A}\left(1-\frac{1}{{N}}\right)+\nu_L+\mu_L\right) \mu_A-\frac{r b \nu_L}{{N}}>0 
 \Leftrightarrow \quad {N}>1.
$$ 
 Since $N>1$ is also an existence condition for $E_{03}$, the analysis confirms that $E_{03}$ is LAS, establishing the stability conditions stated in item 3. On the other hand, $E_{03}$ loses stability when $\frac{r_c}{r_d}<\alpha K_{v_{\max }}\left(1-\frac{1}{N}\right)$ where $\alpha=\frac{m \nu_L}{\mu_A}$.


\item \underline{$E_{04}$:} The Jacobian matrix at $E_{04}$ is given by 
 $$
J=\left[\begin{array}{ccc}
-\left(\frac{r b \nu_L}{\mu_A}\left(1-\frac{1}{{N}}\right)+\nu_L+\mu_L\right) & \frac{r b}{{N}} & -r b \frac{\nu_L}{\mu_A}\left(1-\frac{1}{{N}}\right)^2\left(K_{v_{\max }}-K_{v_{\text {min }}}\right) \\
\nu_L & -\mu_A & 0 \\
0 & 0 & -k\left(-r_c+r_d m\left(\frac{\nu_L}{\mu_A} K_{v_{\text {min }}}\left(1-\frac{1}{{N}}\right)\right)\right.
\end{array}\right].
$$

and the characteristic equation is given by \(\det(J - \lambda I) = 0\) where

\[
%\small{
J-\lambda I = \left[
\begin{array}{ccc}
-\left(\frac{r b \nu_L}{\mu_A}\left(1 - \frac{1}{{N}}\right) + \nu_L + \mu_L\right) - \lambda 
& \frac{r b}{{N}} 
& -r b \frac{\nu_L}{\mu_A}\left(1 - \frac{1}{{N}}\right)^2\left(K_{v_{\max}} - K_{v_{\min}}\right) \\
\nu_L 
& -\mu_A - \lambda 
& 0 \\
0 
& 0 
& -k\left(-r_c + r_d m \left(\frac{\nu_L}{\mu_A} K_{v_{\min}}\left(1 - \frac{1}{{N}}\right)\right)\right) - \lambda
\end{array}
\right].
%}
\]

The determinant of this matrix can be calculated by expanding along the third row, i.e.,

$$
\begin{aligned}
& \operatorname{det}(J-\lambda I)=\left[-k\left(-r_c+r_d m\left(\frac{\nu_L}{\mu_A} K_{v_{\min }}\left(1-\frac{1}{{N}}\right)\right)\right)-\lambda\right] \\
& \operatorname{det}\left[\begin{array}{cc}
-\left(\frac{r b \nu_L}{\mu_A}\left(1-\frac{1}{{N}}\right)+\nu_L+\mu_L\right)-\lambda & \frac{r b}{{N}} \\
\nu_L & -\mu_A-\lambda
\end{array}\right]
\end{aligned}
$$

$$ 
\begin{aligned}
&=\left[-k\left(-r_c+r_d m\left(\frac{\nu_L}{\mu_A} K_{v_{\min }}\left(1-\frac{1}{{N}}\right)\right)\right)-\lambda\right]. \\
& {\left[\left(\frac{r b \nu_L}{\mu_A}\left(1-\frac{1}{{N}}\right)+\nu_L+\mu_L\right)+\lambda\right]\left[\mu_A+\lambda\right]-\left[\frac{r b}{{N}}\right]\left[\nu_L\right] }
\end{aligned}
$$

$$
\begin{aligned}
= & {\left[-k\left(-r_c+r_d m\left(\frac{\nu_L}{\mu_A} K_{v_{\min }}\left(1-\frac{1}{{N}}\right)\right)\right)-\lambda\right] } \\
& {\left[\left(\frac{r b \nu_L}{\mu_A}\left(1-\frac{1}{{N}}\right)+\nu_L+\mu_L\right) \mu_A+\left(\frac{r b \nu_L}{\mu_A}\left(1-\frac{1}{{N}}\right)+\nu_L+\mu_L\right) \lambda+\mu_A \lambda+\lambda^2-\frac{r b \nu_L}{{N}}\right] }
\end{aligned}
$$

$$
\begin{aligned}
& =\left[-k\left(-r_c+r_d t m\left(\frac{\nu_L}{\mu_A} K_{v_{\min }}\left(1-\frac{1}{{N}}\right)\right)\right)-\lambda\right] \\
& {\left[\lambda^2+\left(\left(\frac{r b \nu_L}{\mu_A}\left(1-\frac{1}{{N}}\right)+\nu_L+\mu_L\right)+\mu_A\right) \lambda+\left(\left(\frac{r b \nu_L}{\mu_A}\left(1-\frac{1}{{N}}\right)+\nu_L+\mu_L\right) \mu_A-\frac{r b \nu_L}{{N}}\right)\right]}.
\end{aligned}
$$

Here, eigenvalues: $
\lambda_1=-k\left(-r_c+r_d m\left(\frac{\nu_L}{\mu_A} K_{v_{\min }}\left(1-\frac{1}{{N}}\right)\right)\right) 
$ which is true if $N>1$ and 

$$
\begin{aligned}
& \lambda_{2,3}=\frac{-\left(\left(\frac{r b \nu_L}{\mu_A}\left(1-\frac{1}{{N}}\right)+\nu_L+\mu_L\right)+\mu_A\right)}{2} \pm \\
& \frac{\sqrt{\left(\left(\frac{r b \nu_L}{\mu_A}\left(1-\frac{1}{{N}}\right)+\nu_L+\mu_L\right)+\mu_A\right)^2-4\left(\left(\frac{r b \nu_L}{\mu_A}\left(1-\frac{1}{{N}}\right)+\nu_L+\mu_L\right) \mu_A-\frac{r b \nu_L}{{N}}\right)}}{2}.
\end{aligned}
$$ 



For the steady state to be LAS, the eigenvalue \(\lambda_1\) will be negative if and only if
 $$
 -k\left[-r_c+r_d m\left(\frac{\nu_L}{\mu_A} K_{v_{\min }}\left(1-\frac{1}{{N}}\right)\right)\right]<0 \\
 \Leftrightarrow \frac{r_c}{r_d}<\alpha K_{v_{\min }}\left(1-\frac{1}{{N}}\right) \text {where} \quad \alpha=\frac{m\nu_L}{\mu_A}
$$
and the eigenvalues \(\lambda_{2, 3}\) will have negative real parts if and only if
$$
\left(\left(\frac{r b \nu_L}{\mu_A}\left(1-\frac{1}{{N}}\right)+\nu_L+\mu_L\right) \mu_A-\frac{r b \nu_L}{{N}}\right)>0 
\Leftrightarrow \quad {N}>1.
$$ 
 Since $N>1$ is also an existence condition for $E_{04}$, the analysis confirms that $E_{04}$ is LAS, establishing the stability conditions stated in item 4.
Conversely, $E_{04}$ will be unstable if $\frac{r_c}{r_d}>\alpha K_{v_{\min }}\left(1-\frac{1}{N}\right)$ where $\alpha=\frac{m \nu_L}{\mu_A}$.
\end{itemize}
\end{proof}

\end{document}
